\documentclass[
  spanish,
  letterpaper, oneside
]{article}

\def\documenttitle {Resumen: los derechos humanos y sus garantías en la Constitución Mexicana}
\def\documentsubtitle {Igualdad, libertad, legalidad y seguridad jurídica, y derechos sociales}
\def\documentsubject {Actividad 4 - Garantías Constitucionales}

\def\documentauthor {Martin Jonathan de la Cruz}
\def\coursename {Garantías Constitucionales}
\def\coursecode {025}

\def\universityname {Universidad Abierta y a Distancia de México}
\def\universityfaculty {Licenciatura en Derecho}
\def\universitydepartment {Garantías Constitucionales}
\def\universitydepartmentimage {departamentos/UnADM}
\def\universitydepartmentimagecfg {height=1.57cm}
\def\universitylocation {Roma Norte, Ciudad de México}

\def\authortable {
  \begin{tabular}{ll}
    \textbf{Alumno:}
    & \begin{tabular}[t]{l}
      Martin Jonathan de la Cruz \\
    \end{tabular} \\
    Matrícula: & ES2611202040 \\
    Figura docente: & Aarón López Pérez \\
    Semestre/Bloque: & 2 / 1 \\
    & \\
    \multicolumn{2}{l}{Fecha de realización: \today} \\
    \multicolumn{2}{l}{\universitylocation}
  \end{tabular}
}

\input{template}
\usepackage{helvet}
\renewcommand{\familydefault}{\sfdefault}
\usepackage{array}
\usepackage{longtable}
\usepackage{tikz}
\setcitestyle{authoryear,open={(},close={)}}

\def\coverwatermarkenabled {true}
\def\coverwatermarkimage {img/departamentos/UnADM.pdf}
\def\coverwatermarkopacity {0.16}
\def\coverwatermarkwidth {0.72\paperwidth}
\def\coverwatermarkxshift {0cm}
\def\coverwatermarkyshift {-0.35cm}

\newcommand{\insertcoverwatermark}{%
  \ifthenelse{\equal{\coverwatermarkenabled}{true}}{%
    \AddToShipoutPictureBG*{%
      \begin{tikzpicture}[remember picture,overlay]
        \node[
          opacity=\coverwatermarkopacity,
          inner sep=0pt,
          xshift=\coverwatermarkxshift,
          yshift=\coverwatermarkyshift
        ] at (current page.center) {%
          \includegraphics[width=\coverwatermarkwidth]{\coverwatermarkimage}%
        };
      \end{tikzpicture}%
    }%
  }{}%
}

\begin{document}

\insertcoverwatermark
\templatePortrait
\templatePagecfg
\onehalfspacing

\section{Introducción}

La Constitución Política de los Estados Unidos Mexicanos reconoce los derechos humanos y establece, en el mismo texto, los medios para hacerlos exigibles. Esa doble operación distingue el derecho vigente de una declaración de buenas intenciones: la norma nombra la prerrogativa y, además, fija quién debe respetarla y con qué consecuencia si no lo hace \citep{cpeum2026}.

Este resumen sintetiza el Tema 2 del programa y organiza la parte dogmática constitucional en cuatro grupos: derechos de igualdad, de libertad, de legalidad y seguridad jurídica, y derechos sociales. La clasificación no obedece a una jerarquía entre derechos, pues la Comisión Nacional de los Derechos Humanos advierte que no existen niveles ni jerarquías y que el Estado debe tratarlos de forma global y equitativa \citep{cndhCualesDerechos2026}. Se trata de un criterio expositivo que facilita ubicar cada precepto junto a su garantía correlativa.

Conviene precisar desde el inicio una distinción que atraviesa todo el texto. El \textbf{derecho humano} es la prerrogativa reconocida a la persona; la \textbf{garantía} es el mecanismo jurídico que lo protege y lo hace exigible frente a la autoridad. \textbf{Desde mi perspectiva}, confundirlos vuelve inoperante el análisis: un derecho sin garantía queda como enunciado, y una garantía sin derecho carece de objeto \citep{cndhInehrmConstitucionDH2015}.

\section{Los derechos humanos y sus garantías en la Constitución Mexicana}

\subsection{Fundamento común: el artículo 1.º constitucional}

El artículo 1.º constitucional opera como norma de apertura del sistema. Reconoce que todas las personas gozarán de los derechos humanos previstos en la Constitución y en los tratados internacionales de los que el Estado mexicano sea parte, e incorpora el principio pro persona, que obliga a preferir la interpretación más favorable a la persona \citep{cpeum2026}.

Del mismo precepto derivan cuatro principios rectores, a saber, universalidad, interdependencia, indivisibilidad y progresividad, junto con cuatro obligaciones para toda autoridad: promover, respetar, proteger y garantizar. A ellas se suman los deberes de prevenir, investigar, sancionar y reparar las violaciones a los derechos humanos \citep{cpeum2026}. \textbf{Por tanto}, la reforma constitucional de 2011 desplazó el eje del sistema: la autoridad ya no concede derechos, los reconoce y responde por ellos \citep{cndhInehrmArticulo1o2015}.

\subsection{Derechos humanos de igualdad y sus garantías}

Los derechos de igualdad exigen que la ley trate a las personas sin distinciones arbitrarias. Su expresión más amplia es la prohibición de toda discriminación motivada por origen étnico o nacional, género, edad, discapacidad, condición social o de salud, religión, opiniones, preferencias sexuales o estado civil, contenida en el artículo 1.º, párrafo quinto \citep{cpeum2026}.

La igualdad entre la mujer y el varón ante la ley se establece de forma expresa en el artículo 4.º, párrafo primero, mientras que la composición pluricultural de la nación sustenta el reconocimiento de los pueblos indígenas y su derecho a la libre determinación en el artículo 2.º, apartado A \citep{cpeum2026}. La prohibición de títulos de nobleza y de fueros personales, prevista en el artículo 12, completa el cuadro normativo de este grupo.

La garantía característica opera como una prohibición dirigida a la autoridad, que no puede introducir distinciones lesivas de la dignidad humana. \textbf{En consecuencia}, cuando lo hace, el acto queda sujeto a control constitucional \citep{cndhInehrmConstitucionDH2015}.

\subsection{Derechos humanos de libertad y sus garantías}

Los derechos de libertad delimitan un espacio de autonomía personal donde la autoridad no puede intervenir sin fundamento. La Constitución prohíbe la esclavitud y otorga libertad a quien llegue al territorio nacional en esa condición, según el artículo 2.º, párrafo primero, y reconoce en el artículo 5.º la libertad de trabajo, profesión, industria o comercio, siempre que sean lícitos \citep{cpeum2026}.

En el terreno de las libertades públicas se ubican la manifestación de ideas y el derecho a la información del artículo 6.º, la libertad de imprenta del artículo 7.º, el derecho de petición del artículo 8.º, la libertad de asociación y reunión pacífica del artículo 9.º, el libre tránsito y residencia del artículo 11 y la libertad de convicciones éticas, de conciencia y de religión del artículo 24 \citep{cpeum2026}.

La garantía funciona aquí como límite negativo, pues la autoridad solo puede restringir estas libertades en los casos y bajo las condiciones que la propia Constitución señala. \textbf{Por tanto}, toda restricción no prevista resulta inconstitucional \citep{cndhCualesDerechos2026}.

\subsection{Derechos humanos de legalidad y seguridad jurídica y sus garantías}

Este grupo concentra las garantías procesales que someten al poder público a reglas conocidas de antemano. Su función no consiste en reconocer una prerrogativa material, sino en asegurar que cualquier afectación a la persona ocurra por el cauce debido.

La irretroactividad de la ley en perjuicio de persona alguna y la garantía de audiencia previa al acto privativo se establecen conjuntamente en el artículo 14. El principio de legalidad del artículo 16, párrafo primero, exige que todo acto de molestia conste por escrito, provenga de autoridad competente y funde y motive la causa legal del procedimiento; del mismo precepto derivan el régimen de la detención y la inviolabilidad del domicilio y de las comunicaciones privadas \citep{cpeum2026}.

La prohibición de hacerse justicia por propia mano se acompaña del derecho a la administración de justicia pronta, completa, gratuita e imparcial, conforme al artículo 17. El proceso penal se rige por los principios acusatorio y oral, con presunción de inocencia y defensa adecuada, según el artículo 20, apartado B, y el artículo 22 prohíbe la pena de muerte, la mutilación y los tormentos \citep{cpeum2026}.

\textbf{En este caso}, la garantía tiene naturaleza instrumental, ya que obliga a la autoridad a justificar cada actuación. \textbf{Considero} que este grupo constituye la pieza operativa del sistema: sin fundamentación y motivación exigibles, los demás derechos quedarían sin vía de defensa. La carga de la prueba recae en quien ejerce el poder, no en quien lo resiente \citep{cndhInehrmConstitucionDH2015}.

\subsection{Derechos sociales y sus garantías}

Los derechos sociales imponen al Estado obligaciones de hacer. No basta con que se abstenga de intervenir: debe organizar servicios, destinar recursos y generar las condiciones materiales que hacen posible el ejercicio del derecho.

El artículo 3.º reconoce el derecho a la educación y establece su carácter obligatorio, universal, gratuito y laico. El artículo 4.º concentra el derecho a la alimentación nutritiva, a la protección de la salud, a un medio ambiente sano, al agua y al saneamiento, y a la vivienda digna y decorosa, además del interés superior de la niñez. El derecho al trabajo digno y los derechos de la seguridad social se desarrollan en el artículo 123, apartado A, fracción XXIX \citep{cpeum2026}.

La garantía de estos derechos es progresiva y depende del principio de no regresión, derivado del artículo 1.º, párrafo tercero: el Estado debe avanzar en su satisfacción y no puede retroceder sin justificación estricta \citep{cpeum2026}. \textbf{A diferencia} de los derechos de libertad, cuya exigibilidad es inmediata, los sociales admiten realización gradual conforme a la disponibilidad de recursos \citep{alvarezIcazaDerechosMexico}.

\subsection{El vínculo entre derecho y garantía}

Los cuatro grupos comparten una estructura común, y advertirla evita leerlos como catálogos inconexos. En cada caso concurren tres elementos: el derecho reconocido, la obligación correlativa de la autoridad y el mecanismo de exigibilidad.

% Se enumeran los cuatro grupos porque el tipo de garantía difiere en cada uno;
% la prosa siguiente los reintegra bajo el amparo como cierre del sistema.
\begin{enumerate}
  \item \textbf{Igualdad}: el derecho a no ser discriminado se garantiza mediante la prohibición de distinciones arbitrarias que vincula a toda autoridad (art.~1.º, párr.~5.º CPEUM).
  \item \textbf{Libertad}: el ámbito de autonomía personal se garantiza mediante el principio de reserva, que solo admite las restricciones expresamente previstas (art.~5.º y 6.º CPEUM).
  \item \textbf{Legalidad y seguridad jurídica}: el debido proceso se garantiza mediante la exigencia de fundamentación y motivación de todo acto de autoridad (art.~14 y 16 CPEUM).
  \item \textbf{Derechos sociales}: las condiciones materiales de vida se garantizan mediante obligaciones prestacionales sujetas a progresividad y no regresión (art.~3.º, 4.º y 123 CPEUM).
\end{enumerate}

\textbf{Ahora bien}, cuando la garantía ordinaria resulta insuficiente, el orden constitucional prevé el juicio de amparo como medio de control frente a actos de autoridad que vulneren derechos humanos, conforme a los artículos 103 y 107 \citep{cpeum2026}. Ese mecanismo cierra el sistema, pues convierte el reconocimiento en pretensión procesalmente accionable \citep{marcanoSalazarDerechosHumanos}.

\clearpage
\section{Conclusión}

La revisión del Tema 2 permite sostener una tesis: en el orden constitucional mexicano la clasificación de los derechos humanos tiene valor analítico, no jerárquico. Agrupar los derechos en igualdad, libertad, legalidad y seguridad jurídica, y derechos sociales sirve para identificar qué tipo de obligación genera cada uno y por qué vía se exige, pero no autoriza a subordinar unos a otros, porque el principio de indivisibilidad del artículo 1.º lo impide expresamente \citep{cpeum2026}.

De este ejercicio se sigue una consecuencia práctica para el trabajo jurídico. \textbf{Por tanto}, ante un caso concreto el análisis no termina al nombrar el derecho presuntamente vulnerado: exige precisar el fundamento constitucional con artículo, fracción y párrafo, identificar la obligación correlativa de la autoridad y determinar el mecanismo de exigibilidad aplicable. Un derecho invocado sin su fundamento exacto es un argumento incompleto \citep{cndhInehrmArticulo1o2015}.

\textbf{Sostengo} que el aporte central de la reforma de 2011 no consistió en ampliar el catálogo de derechos, sino en redefinir la posición de la autoridad frente a ellos. Al establecer que todas las autoridades deben promover, respetar, proteger y garantizar los derechos humanos, el artículo 1.º desplazó la carga argumentativa: ya no corresponde a la persona demostrar que merece el derecho, sino a la autoridad justificar cualquier restricción. Esa inversión, más que cualquier enumeración, es lo que vuelve operativo el sistema de garantías.\footnote{En la elaboración de este trabajo se utilizaron herramientas de inteligencia artificial como apoyo para la organización y revisión del texto, conforme a los Lineamientos para el uso ético y responsable de la Inteligencia Artificial en el ámbito académico de la UnADM. El análisis, la selección de fuentes, la verificación de los fundamentos constitucionales y la postura sostenida son de mi autoría.}

\clearpage
\bibliography{garantias-constitucionales}

\end{document}
